\section{Shape-adaptive Fourier cancellation at a prescribed count}
\label{sec:shape-adaptive-fixed-count}
The preceding box theorem assigns one fixed exponent to each coordinate.
Here the frequency region may depend on $n$ without being a box. A
weighted-volume bound replaces the sum of its four width exponents.
We then construct a single nonrectangular region reaching all four
coordinate directions. The logarithmic factor in its definition pays
for the entire dyadic sum; it does not enter the final central rate.

All integrals in this section are in the rescaled variable $v=\sqrt n z$.
The coordinate order remains $(K_1,K_2,N,T)$, with dual physical
coordinates $(u_1,u_2,s,b)$. Put
\begin{equation}\label{eq:shape-fixed-constants}
 \tau_0=\frac1{200},\qquad \theta_0=\frac9{100},\qquad
 s_0=\frac{67}{280},\qquad d_0'=s_0-\tau_0=\frac{41}{175},\qquad
 \eta_0'=\frac14-s_0=\frac3{280}.
\end{equation}
The symbols $d_0'$ and $\eta_0'$ in this section are exponents, not
covariance lower bounds or phase-separation constants.

\subsection{A measurable-region estimate with fixed expansion orders}
\begin{theorem}[A weighted-volume criterion]
\label{thm:shape-volume-criterion}
Fix $C_0<\infty$. For each $n$, let $\Omega_n\subset\R^4$ be any
measurable set such that
\begin{equation}\label{eq:shape-volume-hypotheses}
 \sup_{v\in\Omega_n}|v|\le C_0 n^{\theta_0},\qquad
 |\Omega_n|\le C_0 n^{d_0'},\qquad
 \int_{\Omega_n}|v|\dd v\le C_0 n^{s_0}.
\end{equation}
The supremum may be interpreted essentially. Uniformly in these sets,
$R\in I$, the prescribed count $n$, the mark $0\le k\le n$, and
bounded-BV section insertions $a$, for all sufficiently large $n$,
\begin{align}
 &n^2\int_{\substack{z\in n^{-1/2}\Omega_n\\
                             |z|\ge2n^{-99/200}}}
           |\Phi^{[k],a}_{n,R}(z)|\dd z
 \le C_{C_0}\left(M_a n^{-3/280}+V_a n^{-5559/175}\right).
 \label{eq:shape-volume-cancellation}
\end{align}
No regularity of $\partial\Omega_n$, and no fixed choice of width
exponents for $\Omega_n$, is required.
\end{theorem}
\begin{proof}
Use the exact marked recentering and
\eqref{eq:quarter-characteristic-stopping}. Its integral is at most
$CM_a n^{-1/4}\int_{\Omega_n}|v|\dd v
\le CM_a n^{-3/280}$. The deterministic collision length is
$m=m_k+m_{n-k}=n/c_*+O(1)$ uniformly in $k$.

Apply Theorem~\ref{thm:high-order-damped-unsmoothing} with the fixed choices
\begin{equation}\label{eq:shape-expansion-orders}
 P=26,\qquad Q=29,\qquad
 \delta=\tfrac14n^{-19/100},\qquad
 \epsilon=\tfrac14n^{-32}.
\end{equation}
The residual Taylor degree is $51$ and the spectral degree is $29$.
For $z=v/\sqrt n$ the two analytic margins are $3/100$ and $2/25$,
exactly as for the roof box. The harmless fixed constant $C_0$ is
absorbed by increasing the threshold in $n$. On
$|v|\ge2n^{1/200}$, the damped term is at most a fixed polynomial in
$n$ times $e^{-cn^{1/100}}$, including its $52$ multiplier factors.
Its integral is smaller than every fixed inverse power of $n$.

The fine insertion loss has exponent
$32-d_0'=5559/175$. The fine residual replacement has margin
\begin{equation}\label{eq:shape-fine-margin}
 32-d_0'-51(1/2+\theta_0)=\frac{1173}{700}.
\end{equation}
For the residual term with $b$ nonsingleton blocks, $1\le b\le26$,
\eqref{eq:high-order-damped-unsmoothing} and the volume bound give
\[
 C M_a n^{d_0'+52\theta_0-26+(81/100)b}
                         [\log(2+n)]^{52-b}.
\]
The slowest term has $b=26$ and decay exponent
\begin{equation}\label{eq:shape-residual-margin}
 26(19/100-2\theta_0)-d_0'=\frac9{350}
                 =\frac3{280}+\frac3{200}.
\end{equation}
Every smaller block count improves the power by $81/100$. Thus all
fixed logarithms are absorbed by the strict margin $3/200$.
The $b=1$ contribution is included in this calculation. Finally
$\dd v=n^2\dd z$ and the centering phase has modulus one. This proves
\eqref{eq:shape-volume-cancellation} with constants independent of the
shape of $\Omega_n$. Neither expansion order grows with $n$.
\end{proof}

\subsection{A summable family of coordinate scales}
For $\mathbf j=(j_1,j_2,j_3,j_4)\in\Z_{\ge0}^4$, write
$S(\mathbf j)=\sum_i j_i$ and $J(\mathbf j)=\max_i j_i$.
The following elementary counting fact will be used both for frequency
volume and for the physical convolution kernel.

\begin{lemma}[The weighted dyadic sum]
\label{lem:weighted-dyadic-cross}
For each integer $K\ge0$,
\begin{equation}\label{eq:weighted-dyadic-cross}
 \sum_{S(\mathbf j)+J(\mathbf j)\le K}
                       2^{S(\mathbf j)+J(\mathbf j)}
             \le C2^K(K+1)^3.
\end{equation}
Consequently, if $a>0$, $T\ge a^5$ and $U\ge a$, the finite set
\[
 \mathcal D(a,T,U)=\{\mathbf j:
     a^5 2^{S(\mathbf j)+J(\mathbf j)}\le T,
     \ a2^{J(\mathbf j)}\le U\}
\]
satisfies, with $\lambda_i=a2^{j_i}$ and
$\lambda_*=\max_i\lambda_i$,
\begin{align}
 \sum_{\mathbf j\in\mathcal D}\prod_i\lambda_i\lambda_*
    &\le CT\,[1+\log(T/a^5)]^3,\notag\\
 \sum_{\mathbf j\in\mathcal D}\prod_i\lambda_i
    &\le C(T/a)\,[1+\log(T/a^5)]^3.
 \label{eq:dyadic-cross-scale-sums}
\end{align}
\end{lemma}
\begin{proof}
Fix $t=S+J$ and choose one coordinate attaining the maximum $J$.
For this choice the other three coordinates sum to $t-2J$.
Ignoring their upper bound by $J$ only increases their number, which
is at most $(t+2)^2$. There are at most $t+1$ choices for $J$ and
four choices of its attaining coordinate. Thus the number of tuples
with $S+J=t$ is at most $C(t+1)^3$. Overcounting tied maxima is
harmless. Sum $2^tC(t+1)^3$ for $0\le t\le K$ and use
$\sum_{t=0}^K2^t\le2^{K+1}$. This proves the first assertion.
Take $K=\lfloor\log_2(T/a^5)\rfloor$ and multiply by $a^5$.
The bound on $U$ only removes terms. The second scale sum follows
from $\lambda_*\ge a$.
\end{proof}

Set
\begin{equation}\label{eq:adaptive-cross-scales}
 a_n=2n^{1/200},\quad U_n=n^{9/100},\quad
 T_n=\frac{n^{67/280}}{[\log(2+n)]^3},\quad
 \mathcal D_n=\mathcal D(a_n,T_n,U_n).
\end{equation}
All assertions concerning these sets are for sufficiently large $n$,
so $T_n\ge a_n^5$ and $U_n\ge a_n$. Define the closed union of boxes
\begin{equation}\label{eq:adaptive-cross-support}
 \mathcal H_n=\bigcup_{\mathbf j\in\mathcal D_n}
                       \prod_{i=1}^4[-2\lambda_i,2\lambda_i],
 \qquad \lambda_i=a_n2^{j_i}.
\end{equation}
This is a nonrectangular frequency region. A useful description uses
$x_i(v)=\max\{a_n,|v_i|\}$ and $x_*(v)=\max_i x_i(v)$.
It contains
\begin{equation}\label{eq:adaptive-cross-inner}
 \left\{v:x_*(v)\le U_n/2,\quad
                  x_*(v)\prod_i x_i(v)\le T_n/32\right\},
\end{equation}
and is contained in the same product region with thresholds $2U_n$
and $32T_n$. These constants also allow a smooth cutoff below.

\begin{lemma}[Size and interior of the nonrectangular region]
\label{lem:adaptive-cross-volume}
The set $\mathcal H_n$ satisfies \eqref{eq:shape-volume-hypotheses}
with one fixed constant. It contains $\{|v|\le2n^{1/200}\}$.
The inclusions just stated hold for all sufficiently large $n$.
\end{lemma}
\begin{proof}
Every box has volume $256\prod_i\lambda_i$ and
$|v|\le4\lambda_*$ on it. The two scale sums in
\eqref{eq:dyadic-cross-scale-sums} bound its union volume and first
radial moment. Since $1+\log(T_n/a_n^5)\le C\log(2+n)$, the
factor $[\log(2+n)]^{-3}$ in $T_n$ cancels this entire counting
loss. This gives $|\mathcal H_n|\le Cn^{41/175}$ and
$\int_{\mathcal H_n}|v|\dd v\le Cn^{67/280}$. Its radius is at
most $4U_n$. The zero tuple supplies the central ball.

For the inner inclusion, select scales with
$x_i\le\lambda_i\le2x_i$, taking $j_i=0$ when $x_i=a_n$.
They obey $\lambda_*\prod_i\lambda_i\le32x_*\prod_i x_i\le T_n$
and $\lambda_*\le2x_*\le U_n$. This gives an admissible box
containing $v$. Conversely, membership in a box implies
$x_i\le2\lambda_i$, because $\lambda_i\ge a_n$, and hence the
stated outer bounds.
\end{proof}

\begin{theorem}[One retained-rate region in all four directions]
\label{thm:adaptive-cross-fixed-count}
The physical region is $n^{-1/2}\mathcal H_n$. On its part outside
$|z|\le2n^{-99/200}$,
\begin{align}
 n^2\int |\Phi^{[k],a}_{n,R}(z)|\dd z
       &\le C\left(M_an^{-3/280}+V_an^{-5559/175}\right).
 \label{eq:adaptive-cross-fixed-count}
\end{align}
Let
\begin{equation}\label{eq:adaptive-cross-retained-union}
 \mathcal U_n^\sharp=\mathcal H_n\cup\mathcal B_n^{\rm roof}
                         \cup\{|v|\le2n^{67/1400}\}.
\end{equation}
Then
\begin{align}
 &\int_{\mathcal U_n^\sharp}
 |C^{[k],a}_{n,R}(v)-\alpha_a e^{-v^{\mathsf T}D_Rv/2}|\dd v
 \le C\left(M_an^{-3/280}\sqrt{\log(2+n)}+V_an^{-9/175}\right).
 \label{eq:adaptive-cross-central}
\end{align}
Both estimates are uniform in $R,n,k,a$ in the original marked class.
Every previously proved isotropic and roof-box frequency is retained.
\end{theorem}
\begin{proof}
Theorem~\ref{thm:shape-volume-criterion} and
Lemma~\ref{lem:adaptive-cross-volume} give the first assertion. On
the central ball use Theorem~\ref{thm:marked-return-band}. On its
complement the uniformly positive covariance gives an exponentially
small Gaussian integral, while the first assertion gives the physical
transform integral. Finally add the inherited roof-box and isotropic
bounds. Their overlap can only increase the upper bound. No division
by the volume of any region is made.
\end{proof}

\begin{corollary}[New mixed directions and the unchanged event]
\label{cor:adaptive-cross-directions-event}
For a fixed exponent vector $\boldsymbol\rho$ with $\rho_i\ge0$, set
$\widetilde\rho_i=\max\{1/200,\rho_i\}$ and
$\widetilde\rho_* =\max_i\widetilde\rho_i$. If
\begin{equation}\label{eq:adaptive-cross-exponent-interior}
 \widetilde\rho_*<9/100,\qquad
 \sum_i\widetilde\rho_i+\widetilde\rho_*<67/280,
\end{equation}
then every fixed multiple of the box $|v_i|\le n^{\rho_i}$ is
contained in $\mathcal H_n$ for all sufficiently large $n$.
Also $|v_i|\le U_n/4$ on any one coordinate axis is contained in
$\mathcal H_n$. Thus each physical coordinate has a new axial reach
of order $n^{-41/100}$, not only the roof coordinate.

For the unchanged marked-state event with probability $p_{n,R}>0$
and variation $V_{n,R}$, the right side of
\eqref{eq:adaptive-cross-central} becomes
\[
 \frac C{p_{n,R}}\left(n^{-3/280}\sqrt{\log(2+n)}
                                    +V_{n,R}n^{-9/175}\right).
\]
Its probability and variation conditions are unchanged.
\end{corollary}
\begin{proof}
Apply the inner description \eqref{eq:adaptive-cross-inner}. The
strict power margins absorb both fixed constants and
$[\log(2+n)]^3$. On an axis the product exponent is
$2(9/100)+3/200=39/200<67/280$, so the same argument works at
the asserted constant multiple of $U_n$. As a mixed example,
$\boldsymbol\rho=(1/25,1/25,1/25,1/20)$ has sum plus maximum
$11/50<67/280$ and extends beyond the old isotropic exponent.
The event statement follows by taking $a=c_*^{-1}\one_E$ and
dividing by the exact probability, as in
Corollary~\ref{cor:anisotropic-same-event}.
\end{proof}

\paragraph{Extent of the result.}
This is not the isotropic ball of rescaled radius $n^{1/10}$.
For example, a vector with four coordinates of order $n^{3/50}$
is outside the product budget and outside both retained regions for
large $n$. The expansion orders are fixed. The new result controls a
single explicit nonrectangular domain, rather than a collection of
estimates whose constants might depend on $n$-varying exponents.
The next section connects this domain to the already used raw
extraction without imposing a new, unrelated edge hypothesis.
